% Part: first-order-logic
% Chapter: axiomatic-deduction
% Section: axioms-rules-quantifiers

\documentclass[../../../include/open-logic-section]{subfiles}

\begin{document}

\olfileid{fol}{axd}{qua}

\olsection{પરિમાણકો માટે સ્વયંસિદ્ધિઓ અને નિયમો}

\begin{defn}[પરિમાણકો માટેની સ્વયંસિદ્ધિઓ]
પરિમાણકોને નિયંત્રિત કરતી \emph{સ્વયંસિદ્ધિઓ} નીચેના સ્વરૂપોના બધા
દૃષ્ટાંતો છે:
\begin{align}
\ollabel{ax:q1} & \lforall[x][!B] \lif !B(t), \\
\ollabel{ax:q2} & !B(t) \lif \lexists[x][!B].
\end{align}
જ્યાં $t$ કોઈપણ બંધ પદ છે.
\end{defn}

\begin{defn}[પરિમાણકો માટેના નિયમો]
\begin{enumerate}
\item જો $!B \lif !A(a)$ પહેલેથી !!{derivation}માં આવતું હોય અને
  $a$ એ~$\Gamma$, $!B$ કે~$!A(x)$માં ન આવતું હોય, તો $!B \lif
  \lforall[x][!A(x)]$ યોગ્ય અનુમાન-પગલું છે.
\item જો $!A(a) \lif !B$ પહેલેથી !!{derivation}માં આવતું હોય અને
  $a$ એ~$\Gamma$, $!B$ કે~$!A(x)$માં ન આવતું હોય, તો
  $\lexists[x][!A(x)] \lif
  !B$ યોગ્ય અનુમાન-પગલું છે.
\end{enumerate}
\end{defn}

\sourcecorrection{OLAX-002}{મૂળની
\texttt{axioms-rules-quantifiers.tex}, પંક્તિઓ 23--30માં બંને
\texttt{\string\item} કોઈ યાદી-પર્યાવરણ વિના મૂકવામાં આવ્યાં છે. અહીં
બંને નિયમોને \texttt{enumerate}માં મૂકીને માન્ય લેટેક રચના કરી છે.}

\sourcecorrection{OLAX-003}{મૂળની
\texttt{axioms-rules-quantifiers.tex}, પંક્તિઓ 24--30ની બંને
આઇગનચલ-શરતોમાં $a$ને~$\Gamma$ અને~$!B$માંથી જ બાકાત રાખ્યો છે, પરંતુ
$!A(x)$માંથી નહિ. પછીની યથાર્થતાની સાબિતી માટે પણ આ વધારાની શરત જરૂરી
છે; તેના વિના $!A(x)\ident\eq[x][a]$ લેવાથી બહુ-ઘટક સંરચનામાં અમાન્ય
સર્વત્રીકરણ મળી શકે. તેથી બંને નિયમોમાં $a$ એ~$!A(x)$માં પણ ન આવે એ
શરત ઉમેરી છે.}

આ બંનેમાંથી કોઈપણ નિયમને આપણે ટૂંકમાં ``\QR'' લખીશું.

\end{document}
